% \ifx compares MEANING, not text: two macros with the same replacement text
% are \ifx-equal, and a macro is never \ifx-equal to what it expands to.
\catcode`\{=1 \catcode`\}=2 \catcode`\#=6
\def\one{SAME}
\def\two{SAME}
\def\three{OTHER}
\ifx\one\two \message{one=two}\else \message{one<>two}\fi
\ifx\one\three \message{one=three}\else \message{one<>three}\fi
% An undefined control sequence has a meaning too -- "undefined" -- so two of
% them are \ifx-equal to each other.
\ifx\notdefinedA\notdefinedB \message{undefined=undefined}\fi
% \ifcase selects by number, with \or between the arms and an optional \else
% for everything past the last one.
\count0=2
\message{\ifcase\count0 ZERO\or ONE\or TWO\else MANY\fi}
\count0=9
\message{\ifcase\count0 ZERO\or ONE\or TWO\else MANY\fi}
\end
